\documentclass[12pt, a4paper]{article}

% ─── Packages ────────────────────────────────────────────────────────────────
\usepackage[margin=2.5cm]{geometry}
\usepackage{hyperref}
\usepackage{xcolor}
\usepackage{listings}
\usepackage{titlesec}
\usepackage{enumitem}
\usepackage{fancyhdr}
\usepackage{mdframed}
\usepackage{parskip}
\usepackage{microtype}

% ─── Colours ─────────────────────────────────────────────────────────────────
\definecolor{primary}{HTML}{1D9E75}
\definecolor{secondary}{HTML}{534AB7}
\definecolor{codebg}{HTML}{F5F5F0}
\definecolor{codeborder}{HTML}{D3D1C7}
\definecolor{notebg}{HTML}{E1F5EE}
\definecolor{noteborder}{HTML}{0F6E56}
\definecolor{warnbg}{HTML}{FAEEDA}
\definecolor{warnborder}{HTML}{854F0B}
\definecolor{darktext}{HTML}{2C2C2A}
\definecolor{mutedtext}{HTML}{5F5E5A}
\definecolor{pykeyword}{HTML}{0F6E56}
\definecolor{pycomment}{HTML}{7F7D75}

% ─── Hyperref ────────────────────────────────────────────────────────────────
\hypersetup{
  colorlinks=true,
  linkcolor=secondary,
  urlcolor=primary,
  citecolor=primary,
  pdftitle={LangChain Agent — Quickstart Setup},
  pdfauthor={},
}

% ─── Typography ──────────────────────────────────────────────────────────────
\titleformat{\section}{\Large\bfseries\color{darktext}}{}{0em}{}[\color{primary}\titlerule]
\titleformat{\subsection}{\large\bfseries\color{darktext}}{}{0em}{}

\titlespacing{\section}{0pt}{2em}{0.8em}
\titlespacing{\subsection}{0pt}{1.4em}{0.4em}

% ─── Code listings ───────────────────────────────────────────────────────────
\lstdefinestyle{filestyle}{
  backgroundcolor=\color{codebg},
  basicstyle=\ttfamily\small\color{darktext},
  breaklines=true,
  breakatwhitespace=false,
  frame=single,
  rulecolor=\color{codeborder},
  framesep=8pt,
  xleftmargin=8pt,
  xrightmargin=8pt,
  aboveskip=0.8em,
  belowskip=0.8em,
  showstringspaces=false,
  tabsize=2,
}

\lstdefinestyle{shellstyle}{
  backgroundcolor=\color{codebg},
  basicstyle=\ttfamily\small\color{darktext},
  keywordstyle=\color{pykeyword}\bfseries,
  commentstyle=\color{pycomment}\itshape,
  breaklines=true,
  breakatwhitespace=false,
  frame=single,
  rulecolor=\color{codeborder},
  framesep=8pt,
  xleftmargin=8pt,
  xrightmargin=8pt,
  aboveskip=0.8em,
  belowskip=0.8em,
  showstringspaces=false,
  tabsize=2,
  morekeywords={python,pip,brew,ollama,source,cd,mkdir,git},
}

\lstset{style=filestyle}

% ─── Custom environments ─────────────────────────────────────────────────────
\newmdenv[
  backgroundcolor=notebg,
  linecolor=noteborder,
  linewidth=3pt,
  topline=false,
  bottomline=false,
  rightline=false,
  innerleftmargin=12pt,
  innerrightmargin=12pt,
  innertopmargin=8pt,
  innerbottommargin=8pt,
]{notebox}

\newmdenv[
  backgroundcolor=warnbg,
  linecolor=warnborder,
  linewidth=3pt,
  topline=false,
  bottomline=false,
  rightline=false,
  innerleftmargin=12pt,
  innerrightmargin=12pt,
  innertopmargin=8pt,
  innerbottommargin=8pt,
]{warnbox}

% ─── Header / Footer ─────────────────────────────────────────────────────────
\pagestyle{fancy}
\fancyhf{}
\fancyhead[L]{\small\color{mutedtext}LangChain Agent — Quickstart Setup}
\fancyfoot[C]{\small\color{mutedtext}\thepage}
\renewcommand{\headrulewidth}{0.4pt}
\renewcommand{\headrule}{\hbox to\headwidth{\color{codeborder}\leaders\hrule height \headrulewidth\hfill}}

% ─── Document ─────────────────────────────────────────────────────────────────
\begin{document}

% ── Title page ────────────────────────────────────────────────────────────────
\begin{titlepage}
  \centering
  \vspace*{3cm}
  {\color{primary}\rule{\linewidth}{2pt}}
  \vspace{1cm}

  {\Huge\bfseries\color{darktext} LangChain Agent}\\[0.5em]
  {\Large\color{mutedtext} Quickstart Setup Guide}

  \vspace{1cm}
  {\color{primary}\rule{\linewidth}{2pt}}

  \vspace{2cm}
  {\large\color{mutedtext}
    Everything you need to install and run\\
    \texttt{1-langchain-agent/agent.py} from scratch.
  }

  \vfill
  {\small\color{mutedtext} Version 1.0}
\end{titlepage}

\tableofcontents
\newpage

% ══════════════════════════════════════════════════════════════════════════════
\section{Step 1 — Install Ollama}
% ══════════════════════════════════════════════════════════════════════════════

Ollama runs large language models locally. Install it from \url{https://ollama.com/download} or via Homebrew on macOS:

\begin{lstlisting}[style=shellstyle]
# macOS with Homebrew
brew install ollama
\end{lstlisting}

For Linux, use the official install script:

\begin{lstlisting}[style=shellstyle]
curl -fsSL https://ollama.com/install.sh | sh
\end{lstlisting}

% ══════════════════════════════════════════════════════════════════════════════
\section{Step 2 — Start the Ollama Server}
% ══════════════════════════════════════════════════════════════════════════════

Open a \textbf{separate terminal} and start the Ollama daemon. Keep this terminal open for the duration of your session:

\begin{lstlisting}[style=shellstyle]
ollama serve
\end{lstlisting}

\begin{notebox}
On macOS, launching the Ollama desktop app starts the server automatically. If you installed via Homebrew and prefer the CLI, \texttt{ollama serve} is the equivalent.
\end{notebox}

% ══════════════════════════════════════════════════════════════════════════════
\section{Step 3 — Pull the Model}
% ══════════════════════════════════════════════════════════════════════════════

In your original terminal, pull the model the agent is configured to use:

\begin{lstlisting}[style=shellstyle]
ollama pull gemma4:e4b
\end{lstlisting}

This downloads the model weights to your local machine. The download is a one-time operation.

\begin{notebox}
If you prefer a different local model, pull it with \texttt{ollama pull <name>} and update the \texttt{model} argument in \texttt{agent.py} accordingly. The model must support the chat/instruct format; base models will not follow the pipeline instructions reliably.
\end{notebox}

% ══════════════════════════════════════════════════════════════════════════════
\section{Step 4 — Clone the Repository}
% ══════════════════════════════════════════════════════════════════════════════

\begin{lstlisting}[style=shellstyle]
git clone https://github.com/ProgettoLabs/report-agent.git
cd report-agent
\end{lstlisting}

% ══════════════════════════════════════════════════════════════════════════════
\section{Step 5 — Create and Activate a Virtual Environment}
% ══════════════════════════════════════════════════════════════════════════════

\begin{lstlisting}[style=shellstyle]
python3 -m venv venv
source venv/bin/activate
\end{lstlisting}

On Windows (PowerShell):

\begin{lstlisting}[style=shellstyle]
python -m venv venv
venv\Scripts\Activate.ps1
\end{lstlisting}

% ══════════════════════════════════════════════════════════════════════════════
\section{Step 6 — Install Python Dependencies}
% ══════════════════════════════════════════════════════════════════════════════

The simple agent only needs two packages:

\begin{lstlisting}[style=shellstyle]
pip install langchain-core langchain-ollama
\end{lstlisting}

No other dependencies are required to run \texttt{1-langchain-agent/agent.py}.

% ══════════════════════════════════════════════════════════════════════════════
\section{Step 7 — Create a Pipeline Directory}
% ══════════════════════════════════════════════════════════════════════════════

The agent expects a specific directory layout. Create your pipeline folder anywhere inside the repository (the \texttt{use-cases/} directory is conventional):

\begin{lstlisting}[language={}]
use-cases/
+-- my-pipeline/
    |-- pipeline.md          # High-level description of what the pipeline does
    |-- input_data.md        # The raw data the pipeline will process
    |-- step_01_<name>/
    |   |-- spec.md          # What this step must do
    |   +-- output.md        # The required output format
    |-- step_02_<name>/
    |   |-- spec.md
    |   +-- output.md
    +-- step_03_<name>/
        |-- spec.md
        +-- output.md
\end{lstlisting}

\begin{itemize}[leftmargin=1.5em]
  \item \textbf{\texttt{pipeline.md}} — a short description of the overall pipeline goal, shown to the model at every step.
  \item \textbf{\texttt{input\_data.md}} — the raw data fed into the pipeline (CSV, JSON, prose, or any text format).
  \item \textbf{\texttt{step\_NN\_<name>/}} — one directory per step, sorted lexicographically. The zero-padded numeric prefix (\texttt{01}, \texttt{02}, \ldots) determines execution order.
  \item \textbf{\texttt{spec.md}} — instructions for the LLM describing exactly what the step should do.
  \item \textbf{\texttt{output.md}} — the exact format the LLM must produce (with an example if helpful).
\end{itemize}

\begin{notebox}
A working example pipeline is included in the repository at \texttt{use-cases/1-quarterly-report/}. Inspect it to see how each file is written before creating your own.
\end{notebox}

% ══════════════════════════════════════════════════════════════════════════════
\section{Step 8 — Run the Agent}
% ══════════════════════════════════════════════════════════════════════════════

Pass the pipeline directory as the sole argument:

\begin{lstlisting}[style=shellstyle]
python 1-langchain-agent/agent.py use-cases/my-pipeline
\end{lstlisting}

The agent will print progress for each step. When it finishes, two files will have been written inside the pipeline directory:

\begin{itemize}[leftmargin=1.5em]
  \item \textbf{\texttt{context.json}} — the full output of every step, keyed by step directory name.
  \item \textbf{\texttt{final\_report.md}} — a verbatim copy of the last step's output.
\end{itemize}

\begin{warnbox}
Make sure the Ollama server from Step 2 is still running before executing the agent. If the connection is refused, restart \texttt{ollama serve} in its terminal.
\end{warnbox}

\end{document}
